\section{Background: How Retrieval-Augmented Generation Works}
\label{sec:background}

This section is a self-contained primer. Readers familiar with dense
retrieval and RAG can skip to Section~\ref{sec:related}; every concept
introduced here is used later to explain a measured result.

\subsection{The knowledge problem, and why not simply retrain the model}
\label{sec:why-rag}

A large language model stores what it knows in its
weights~\cite{brown2020gpt3,petroni2019lmkb} --- so-called
\emph{parametric} knowledge, acquired once, during training. Three
consequences follow. First, the knowledge is \emph{frozen}: a model trained
in 2024 knows nothing that happened afterwards. Second, it is
\emph{incomplete and unreliable}: how much a model retains depends on how
often it saw a fact~\cite{roberts2020knowledge,kandpal2023longtail} ---
popular facts appear thousands of times where rare facts appear once or
never~\cite{mallen2023popqa} --- and when asked about the latter it tends
to produce fluent, plausible, wrong answers ---
hallucinations~\cite{ji2023hallucination,zhang2023sirens}. Third, it is
\emph{closed}: an
organisation's internal documentation, product catalogue, or case history
was never in the training data at all.

The obvious remedy --- continue training the model on the missing documents
--- is rarely practical. Full retraining of even a mid-sized model costs
GPU-months; parameter-efficient fine-tuning~\cite{hu2022lora} is cheaper
but must be repeated every time the documents change, offers no
attribution (the model still cannot say \emph{where} an answer came
from), and injects new factual knowledge less reliably than simply
retrieving it~\cite{ovadia2024finetuning}. Retrieval-Augmented
Generation~\cite{lewis2020rag} sidesteps the problem: the documents stay
outside the model, in a searchable index that can be updated by inserting
or deleting a record, and at question time the system retrieves the few
passages most relevant to the question and prepends them to the model's
prompt. The model's job shrinks from \emph{recalling} the fact to
\emph{reading} it. Updating knowledge becomes a database operation rather than
a training run; a company can place its private business context into the
index without any of it entering the model's weights; and every answer can
cite the passages it was read from.

\subsection{Embeddings and cosine similarity}
\label{sec:embeddings}

The search step must match a question to passages by \emph{meaning}, not
just by shared words --- ``Where did the director of Beat Girl die?'' shares
almost no vocabulary with the encyclopedia paragraph that answers it. The
standard tool is the \emph{text embedding}: a neural network (the
\emph{embedding model} or \emph{encoder}, typically a Transformer
encoder~\cite{vaswani2017attention,devlin2019bert}) that maps a piece of
text to a vector of a few hundred real numbers, trained so that texts
with similar meaning land close together in that vector
space~\cite{reimers2019sbert}.

Closeness is measured by \emph{cosine similarity}: the cosine of the angle
between two vectors, which is $1$ for parallel vectors, $0$ for orthogonal
ones, and $-1$ for opposite ones. For vectors $\mathbf{a}$ and $\mathbf{b}$,
\begin{equation}
\cos(\mathbf{a},\mathbf{b})
  = \frac{\mathbf{a}\cdot\mathbf{b}}{\lVert\mathbf{a}\rVert\,\lVert\mathbf{b}\rVert}
  = \frac{\sum_i a_i b_i}{\sqrt{\sum_i a_i^2}\sqrt{\sum_i b_i^2}}.
\end{equation}
A toy example with three dimensions makes the mechanics concrete. Suppose an
encoder mapped three texts as follows:
\begin{center}
\small
\begin{tabular}{lr}
\toprule
Text & Vector \\
\midrule
``Who directed Beat Girl?'' & $(0.9,\; 0.4,\; 0.1)$ \\
``Beat Girl is a 1960 film directed by Edmond T.\ Gr\'eville.'' & $(0.8,\; 0.5,\; 0.2)$ \\
``The recipe calls for two eggs and a cup of flour.'' & $(0.1,\; 0.2,\; 0.9)$ \\
\bottomrule
\end{tabular}
\end{center}
The question and the film passage give
$\cos = (0.72+0.20+0.02)/(1.0\cdot0.97) \approx 0.97$; the question and the
recipe give $\cos \approx 0.35$. The film passage is retrieved. Real
encoders behave the same way in 768 dimensions instead of three.

Two properties of embeddings matter for the results in this article. First,
an embedding is a \emph{lossy compression}: 768 numbers cannot preserve
everything in a hundred-word passage, and what the compression loses first
is rare surface forms --- exact titles, names, codes. A question that names
its target verbatim can therefore still miss it in embedding space, which
is precisely the failure that lexical search repairs
(Section~\ref{sec:arch-hybrid}). Second, some encoders are trained
\emph{asymmetrically}: the query is expected to carry a short instruction
prefix (e.g.\ \texttt{query:}~\ldots) that the passage side does not.
Ignoring the convention silently degrades retrieval; we verify the
convention for our encoder empirically (Section~\ref{sec:setup}).

\subsection{Embedding models and cross-encoders}
\label{sec:encoders}

The encoder used for search is a \emph{bi-encoder}: it embeds the question
and each passage \emph{separately}, and their relevance is the similarity
of the two vectors. This is what makes web-scale search possible --- all
passage vectors are computed once, offline, and a question at query time is
embedded once and compared against millions of stored vectors in
milliseconds. The price is that the encoder never sees the question and the
passage \emph{together}, so it cannot model their interaction; it matches
what a question is \emph{about}.

A \emph{cross-encoder}~\cite{nogueira2019passage} reads the concatenated
(question, passage) pair through one network and outputs a relevance score.
Seeing both texts jointly, it can judge whether the passage actually
answers the question --- what the question is \emph{asking for} --- far
more precisely, but at the cost of one full network pass per pair, which
rules out scanning a corpus. The standard compromise is a two-stage
pipeline: a bi-encoder retrieves a shortlist of candidates, and the
cross-encoder reorders the shortlist (Section~\ref{sec:arch-rerank}). The
distinction between \emph{about} and \emph{asking for} is not cosmetic; it
produces the single largest failure we measure
(Section~\ref{sec:results-single}).

\subsection{Vector databases}
\label{sec:vectordb}

Storing millions of embedding vectors and answering ``which stored vectors
are most similar to this query vector'' is the job of a \emph{vector
database}. Comparing the query against every stored vector (exact search)
is linear in corpus size; even GPU-accelerated exact
search~\cite{johnson2019faiss} is a poor fit for serving interactive
queries at the 21-million-passage scale used here. Vector databases therefore build an
\emph{approximate nearest neighbour} (ANN) index --- most commonly a
Hierarchical Navigable Small World graph (HNSW)~\cite{malkov2020hnsw}, a
layered proximity graph walked greedily from an entry point towards the
query. The walk inspects a small fraction of the corpus and returns the
nearest vectors it visited, trading a tunable amount of recall for orders of
magnitude in speed; the size of the candidate list the walk keeps
(\texttt{ef}) controls that trade, and we report and fix it in our setup
because a benchmark that leaves it at an unknown default silently measures
the index rather than the architecture.

Alongside each vector, the database stores the passage text and its
identifier, so a search returns readable passages, ready to be placed in a
prompt. The same database can also host a \emph{sparse} (lexical) index, in
which a passage is represented by weighted term counts and matching is by
exact term overlap under BM25 scoring~\cite{robertson2009bm25} --- the
modern packaging of classical keyword search, and the complement to dense
retrieval used by the hybrid architectures below.

\subsection{The naive RAG pipeline}
\label{sec:naive-pipeline}

\begin{figure}[t]
\centering
\begin{tikzpicture}[node distance=5mm and 8mm]
% Ingestion phase
\node[doc] (corpus) {Documents};
\node[comp, fill=black!5, right=of corpus] (chunk) {Split into\\$\sim$100-word passages};
\node[ret, right=of chunk] (encdoc) {Embedding\\model};
\node[store, right=of encdoc] (db) {Vector\\database};
\draw[arr] (corpus) -- (chunk);
\draw[arr] (chunk) -- (encdoc);
\draw[arr] (encdoc) -- node[lbl, above] {vectors} (db);
\node[lbl, anchor=west] at ($(corpus.west)+(-2mm,9mm)$) {\textbf{Ingestion (once, offline)}};
% Query phase
\node[q, below=13mm of corpus] (question) {Question};
\node[ret, right=of question] (encq) {Embedding\\model};
\node[comp, fill=blue!6, right=of encq] (search) {Similarity\\search};
\node[llm, right=of search] (gen) {Generator\\(LLM)};
\node[ans, right=of gen] (answer) {Answer};
\draw[arr] (question) -- (encq);
\draw[arr] (encq) -- node[lbl, above] {vector} (search);
\draw[arr] (db.south) to[out=-90,in=90] node[lbl, right] {top-$k$ passages} (search.north);
\draw[arr] (search) -- (gen);
\draw[arr] (gen) -- (answer);
\draw[arr] (question.south) to[out=-60,in=-150] ($(gen.south)+(-2mm,0)$);
\node[lbl, anchor=west] at ($(question.west)+(-2mm,9mm)$) {\textbf{Query time (per question)}};
\end{tikzpicture}
\caption{The naive RAG pipeline. Offline, the document collection is split
into passages, embedded, and stored in a vector database. At query time the
question is embedded with the same model, the $k$ most similar passages are
retrieved, and the generator answers from the question plus those passages.
Every architecture in this article is a modification of one or more blocks
of this pipeline.}
\label{fig:naive-pipeline}
\end{figure}

Figure~\ref{fig:naive-pipeline} assembles the components into the baseline
architecture that every refinement modifies. Offline, the collection is
split into passages of roughly one hundred words --- small enough that a
retrieved passage is mostly about one thing, large enough to carry a fact
in context --- and each passage is embedded and stored. At query time, the
question is embedded, the $k$ most similar passages are retrieved (here
$k=5$), and the generator receives a prompt of the form ``Answer the
question based on the given documents\ldots'' followed by the passages and
the question.

Everything that can go wrong in this pipeline goes wrong in one of a few
ways: the passage that answers the question is in the corpus but ranked
below the top~$k$; it is unreachable because the question's own words are
too far from it in embedding space; the right article was found but the
wrong hundred-word slice of it was taken; or the answer simply is not in
the corpus. Which of these dominates depends on the question distribution,
and measuring that split --- rather than guessing it --- is what
Section~\ref{sec:results-diagnosis} contributes; the architectures of
Section~\ref{sec:architectures} are, in effect, one remedy per failure
mode. The generation side has failure modes of its own --- readers use
passages unevenly by position in the context~\cite{liu2024lost}, react
measurably to irrelevant passages~\cite{cuconasu2024noise}, and
engineering surveys catalogue further failure points across the
pipeline~\cite{barnett2024seven} --- but on this benchmark the bulk of
the measurable headroom sits in retrieval, which is where the diagnosis
looks.
